\scanpage{161}
\mspage{3.}
\margin{160}
Les sites annelés $\mathrm{Ouv}(X)$ et $\mathrm{Ouv}(Y)$ sont donc isomorphes. Ils ont donc la m\inserted{ê}me cohomologie et, de plus, leurs faisceaux (quasi-) cohérents, localement libres etc. se correspondent, d'où la relation avec la théorie classique. Pour plus de détails voir \inserted{[}EGA IV , § 10\inserted{]} déjà cité.

Nous rencontrons déjà ici un fait qui va \inserted{ê}tre illustré plus abondamment dans ce qui suit : ``dans la pratique'' un site n'est pas donné d'une façon intrinsèque. Ici les sites $\mathrm{Ouv}(X)$ et $\mathrm{Ouv}(Y)$ sont cependant équivalentes (et m\inserted{ê}me isomorphes ). On verra que la catégorie des faisceaux associée à un site (i.e. le topos associé ) est beaucoup plus intrinsèque. Choisir un site particulier admettant un topos donné comme topos associé est une opération du m\inserted{ê}me genre que de choisir un système de générateurs d'un espace vectoriel, et il est naturel de chercher un tel site qui est aussi ``petit'' que possible.

On peut cependant introduire à l'aide des sites $\mathrm{Ouv}(Z)$ ($Z/X\in(\mathrm{Sch})_{/X}$) une topologie bien déterminée sur $(\mathrm{Sch})_{/X}$ de la manière suivante:

L'application qui a chaque ouvert de Zariski d'un préschéma $Z$ associe le préschéma induit sur cet ouvert définit un foncteur
\[
\mathrm{Ouv}(Z)\xrightarrow{i_Z}(\mathrm{Sch})_{/Z}
\]
d'où un foncteur sur les catégories des préfaisceaux
\[
\widehat{\mathrm{Ouv}(Z)}\xleftarrow{i_{Z\bullet}}\widehat{(\mathrm{Sch})_{/Z}}
\]

\uline{Proposition et définition 1.}-\uline{La sous-catégorie} (strictement) \uline{pleine de} $\widehat{(\mathrm{Sch})_{/X}}$ \uline{formée des préfaisceaux} $F$ \uline{tels que} $i_{Z\bullet}(F)\in\widetilde{\mathrm{Ouv}(Z)}$ \uline{pour tout} $Z/X\in(\mathrm{Sch})_{/X}$ \uline{est une catégorie de} $\underline{U}$-\uline{faisceaux sur} $(\mathrm{Sch})_{/X}$. \deleted{\illegible{5}} \uline{La topologie correspondante sur} $(\mathrm{Sch})_{/X}$ \uline{peut \inserted{ê}tre engendrée par la prétopologie où une famille} $\{Z_\alpha\xrightarrow{j_\alpha}Z\}_{\alpha\in I}$ \uline{de morphismes de} $(\mathrm{Sch})_{/X}$ \uline{est une famille couvrante si et seulement si}
